\chapter{Introduction}\label{ch:introduction}

In 11 of the 23 years between 2001 and 2023, the loss ratio of Korea's crop disaster insurance exceeded 100 percent: claims were larger than the premiums set aside to cover them. Losses of this kind are hard for private insurers to carry, because weather shocks strike many farms at once and insurers observe farm risk and effort only imperfectly \citep{miranda1997,mahul2010}. South Korea responds with premium subsidies, administrative cost subsidies, and state reinsurance. In 2026, the program covered 78 crops, with KRW 556.6 billion budgeted for national premium and loading support (Ministry of Agriculture, Food and Rural Affairs [MAFRA], \citeyear{mafra2026}).

A rationale for intervention does not settle its scale or form. Premium support changes the price of protection, administrative support finances delivery, and reinsurance reallocates losses, so each instrument answers a different question. Korean studies have focused instead on how production, input use, or income differs for insured crops or farms \citep{han2014,nam2022,leekim2026}. They do not evaluate the instruments separately, and production, only one margin of performance, says little directly about income protection, distributional benefits, or fiscal efficiency.

This thesis asks two questions: how do premium subsidies, administrative cost subsidies, and state reinsurance address constraints in Korea's crop insurance market, and what production patterns accompany insurance availability? The first is addressed by assessing each instrument's financing and incentives using program rules and financing data. The second is addressed by comparing seven annual crops with crops not yet exposed to insurance, and early with later rice-pilot municipalities. Separating area from yield, and varying reference periods, comparison crops, and weights, shows which choices drive the production contrasts.

The central argument is that each instrument should be judged against its own objective rather than against aggregate production. The institutional assessment points to three priorities: evaluating premium transfers against access and distributional goals, linking administrative support to cost and service quality, and comparing reinsurance arrangements by how they allocate and price risk. Production evidence alone cannot carry this judgment, a limitation that applies, to be fair, to the present comparisons: the national contrasts are imprecise and sensitive to comparison choices, and the rice comparison, although it holds crop identity fixed, covers only the first two years of availability. The analysis concerns insurance availability as a policy package, not the separate effect of each instrument or the welfare of enrolled farms. Chapters 2 and 3 describe the program and assess government support, Chapters 4 and 5 present the data, methods, and findings, and Chapters 6 and 7 discuss policy implications and conclude.
